% !TEX program = tectonic
% =============================================================================
%  AULA 6 — Redes sem Fio: Calidade de Servizo (QoS)
%  Curso: Redes sem Fio  |  Profa. Anelise Munaretto
%  Intserv/Diffserv, RSVP, WFQ, token bucket, 802.11e, QoS ad hoc (QOLSR).
%  Compilar: tectonic aula6.tex
% =============================================================================
\documentclass[aspectratio=169,11pt]{beamer}
% =============================================================================
%  PREÁMBULO COMÚN para as aulas de Redes sem Fio (Beamer + TikZ)
%  Incluido com % =============================================================================
%  PREÁMBULO COMÚN para as aulas de Redes sem Fio (Beamer + TikZ)
%  Incluido com % =============================================================================
%  PREÁMBULO COMÚN para as aulas de Redes sem Fio (Beamer + TikZ)
%  Incluido com \input{preamble.tex} en cada aula.
%  Paleta de cores e estilos são definidos aqui uma única vez.
% =============================================================================

\usepackage[T1]{fontenc}
\usepackage{amsmath,amssymb,mathtools}
\usepackage{graphicx}
\usepackage{tikz}
\usetikzlibrary{positioning,arrows.meta,calc,backgrounds,decorations.pathmorphing,fit,shadows,shapes.symbols,shapes.geometric}
\usepackage{pgfplots}
\pgfplotsset{compat=1.16}
\usepackage{tabularx,booktabs,multirow,array}
\usepackage[normalem]{ulem}
\usepackage{hyperref}

% ---------- Paleta de cores própria ----------
\definecolor{aneliseAzul}{RGB}{20,60,110}
\definecolor{aneliseVerde}{RGB}{30,140,70}
\definecolor{aneliseLaranja}{RGB}{215,120,20}
\definecolor{aneliseGris}{RGB}{90,90,90}

% ---------- Tema ----------
\usetheme{Madrid}
\usecolortheme{whale}
\setbeamercolor{structure}{fg=aneliseAzul}
\setbeamercolor{title}{fg=aneliseAzul}
\setbeamercolor{frametitle}{fg=aneliseAzul}
\setbeamercolor{block title}{fg=aneliseAzul,bg=aneliseGris!12!white}
\setbeamercolor{block body}{bg=aneliseGris!7!white}
\hypersetup{colorlinks=true,linkcolor=aneliseAzul,urlcolor=aneliseVerde}

% ---------- Comandos auxiliares ----------
\newcommand{\kurose}{Kurose \& Ross, \emph{Computer Networking: a Top-Down Approach}}
\newcommand{\forouzan}{Forouzan, \emph{Data Communications and Networking}}
\newcommand{\stallings}{Stallings, \emph{Wireless Communications \& Networks}}
\newcommand{\tanenbaum}{Tanenbaum, \emph{Computer Networks}}
\newcommand{\rfc}[1]{\href{https://www.rfc-editor.org/rfc/rfc#1.txt}{RFC~#1}}
% -----------------------------------------------------------------------------
%  Estilos TikZ comuns para desenhar redes (posicionamento relativo)
% -----------------------------------------------------------------------------
\tikzset{
  host/.style={draw,rounded corners=2mm,fill=aneliseAzul!15,inner sep=1.5mm,font=\scriptsize},
  rt/.style={draw,rounded corners=1mm,fill=aneliseLaranja!25,inner sep=1mm,font=\scriptsize},
  nube/.style={draw,cloud,aspect=2.2,fill=aneliseGris!15,inner sep=1.5mm,font=\scriptsize},
  el/.style={draw,rounded corners=1.2mm,fill=aneliseGris!18,inner sep=0.8mm,font=\scriptsize},
  enl/.style={thick},
  enlA/.style={thick,-Stealth},
} en cada aula.
%  Paleta de cores e estilos são definidos aqui uma única vez.
% =============================================================================

\usepackage[T1]{fontenc}
\usepackage{amsmath,amssymb,mathtools}
\usepackage{graphicx}
\usepackage{tikz}
\usetikzlibrary{positioning,arrows.meta,calc,backgrounds,decorations.pathmorphing,fit,shadows,shapes.symbols,shapes.geometric}
\usepackage{pgfplots}
\pgfplotsset{compat=1.16}
\usepackage{tabularx,booktabs,multirow,array}
\usepackage[normalem]{ulem}
\usepackage{hyperref}

% ---------- Paleta de cores própria ----------
\definecolor{aneliseAzul}{RGB}{20,60,110}
\definecolor{aneliseVerde}{RGB}{30,140,70}
\definecolor{aneliseLaranja}{RGB}{215,120,20}
\definecolor{aneliseGris}{RGB}{90,90,90}

% ---------- Tema ----------
\usetheme{Madrid}
\usecolortheme{whale}
\setbeamercolor{structure}{fg=aneliseAzul}
\setbeamercolor{title}{fg=aneliseAzul}
\setbeamercolor{frametitle}{fg=aneliseAzul}
\setbeamercolor{block title}{fg=aneliseAzul,bg=aneliseGris!12!white}
\setbeamercolor{block body}{bg=aneliseGris!7!white}
\hypersetup{colorlinks=true,linkcolor=aneliseAzul,urlcolor=aneliseVerde}

% ---------- Comandos auxiliares ----------
\newcommand{\kurose}{Kurose \& Ross, \emph{Computer Networking: a Top-Down Approach}}
\newcommand{\forouzan}{Forouzan, \emph{Data Communications and Networking}}
\newcommand{\stallings}{Stallings, \emph{Wireless Communications \& Networks}}
\newcommand{\tanenbaum}{Tanenbaum, \emph{Computer Networks}}
\newcommand{\rfc}[1]{\href{https://www.rfc-editor.org/rfc/rfc#1.txt}{RFC~#1}}
% -----------------------------------------------------------------------------
%  Estilos TikZ comuns para desenhar redes (posicionamento relativo)
% -----------------------------------------------------------------------------
\tikzset{
  host/.style={draw,rounded corners=2mm,fill=aneliseAzul!15,inner sep=1.5mm,font=\scriptsize},
  rt/.style={draw,rounded corners=1mm,fill=aneliseLaranja!25,inner sep=1mm,font=\scriptsize},
  nube/.style={draw,cloud,aspect=2.2,fill=aneliseGris!15,inner sep=1.5mm,font=\scriptsize},
  el/.style={draw,rounded corners=1.2mm,fill=aneliseGris!18,inner sep=0.8mm,font=\scriptsize},
  enl/.style={thick},
  enlA/.style={thick,-Stealth},
} en cada aula.
%  Paleta de cores e estilos são definidos aqui uma única vez.
% =============================================================================

\usepackage[T1]{fontenc}
\usepackage{amsmath,amssymb,mathtools}
\usepackage{graphicx}
\usepackage{tikz}
\usetikzlibrary{positioning,arrows.meta,calc,backgrounds,decorations.pathmorphing,fit,shadows,shapes.symbols,shapes.geometric}
\usepackage{pgfplots}
\pgfplotsset{compat=1.16}
\usepackage{tabularx,booktabs,multirow,array}
\usepackage[normalem]{ulem}
\usepackage{hyperref}

% ---------- Paleta de cores própria ----------
\definecolor{aneliseAzul}{RGB}{20,60,110}
\definecolor{aneliseVerde}{RGB}{30,140,70}
\definecolor{aneliseLaranja}{RGB}{215,120,20}
\definecolor{aneliseGris}{RGB}{90,90,90}

% ---------- Tema ----------
\usetheme{Madrid}
\usecolortheme{whale}
\setbeamercolor{structure}{fg=aneliseAzul}
\setbeamercolor{title}{fg=aneliseAzul}
\setbeamercolor{frametitle}{fg=aneliseAzul}
\setbeamercolor{block title}{fg=aneliseAzul,bg=aneliseGris!12!white}
\setbeamercolor{block body}{bg=aneliseGris!7!white}
\hypersetup{colorlinks=true,linkcolor=aneliseAzul,urlcolor=aneliseVerde}

% ---------- Comandos auxiliares ----------
\newcommand{\kurose}{Kurose \& Ross, \emph{Computer Networking: a Top-Down Approach}}
\newcommand{\forouzan}{Forouzan, \emph{Data Communications and Networking}}
\newcommand{\stallings}{Stallings, \emph{Wireless Communications \& Networks}}
\newcommand{\tanenbaum}{Tanenbaum, \emph{Computer Networks}}
\newcommand{\rfc}[1]{\href{https://www.rfc-editor.org/rfc/rfc#1.txt}{RFC~#1}}
% -----------------------------------------------------------------------------
%  Estilos TikZ comuns para desenhar redes (posicionamento relativo)
% -----------------------------------------------------------------------------
\tikzset{
  host/.style={draw,rounded corners=2mm,fill=aneliseAzul!15,inner sep=1.5mm,font=\scriptsize},
  rt/.style={draw,rounded corners=1mm,fill=aneliseLaranja!25,inner sep=1mm,font=\scriptsize},
  nube/.style={draw,cloud,aspect=2.2,fill=aneliseGris!15,inner sep=1.5mm,font=\scriptsize},
  el/.style={draw,rounded corners=1.2mm,fill=aneliseGris!18,inner sep=0.8mm,font=\scriptsize},
  enl/.style={thick},
  enlA/.style={thick,-Stealth},
}

\title[Redes sem Fio]{Redes sem Fio\\ Calidade de Servizo (QoS)}
\subtitle{Aula 6}
\author[Anelise Munaretto]{Profa. Anelise Munaretto}
\institute[Curso]{Redes sem Fio --- Ciência da Computación}
\date{2026}

\begin{document}

\begin{frame}[plain]
  \titlepage
  \begin{center}
    \scriptsize Soporte QoS en redes IP: Intserv, Diffserv, RSVP; QoS en WiFi (802.11e) e ad hoc.
  \end{center}
\end{frame}

\section{Multimedia e QoS}

\begin{frame}{Multimídia e calidade de servizo}
  \begin{columns}
    \begin{column}{0.5\textwidth}
      \textbf{Aplicacións de multimídia}
      \begin{itemize}
        \item audio e vídeo de rede (``mídia continua'')
        \item requiren atraso e perda controlados
      \end{itemize}
    \end{column}
    \begin{column}{0.5\textwidth}
      \textbf{QoS}
      \begin{itemize}
        \item a rede ofrece á aplicación o \textbf{nivel de desempeño necesario} para que funcione
      \end{itemize}
    \end{column}
  \end{columns}
  \vspace{0.3cm}
  \begin{alertblock}{Problema}
    TCP/UDP/IP: ``servizo de mellor esforzo'' --- \textbf{sen garantía sobre atraso e perda}.
  \end{alertblock}
  \begin{block}{Estratexias}
    \begin{itemize}
      \item aplicacións usan técnicas de capa de aplicación para aliviar atraso/perda (buffering, adaptación de taxa, FEC)
      \item a rede debe evolucionar para dar mellor soporte (clases de servicio)
    \end{itemize}
  \end{block}
\end{frame}

\begin{frame}{Como debe evolucionar a Internet para multimedia?}
  \begin{columns}
    \begin{column}{0.5\textwidth}
      \textbf{Filosofía de servizos integrados (Intserv)}
      \begin{itemize}
        \item cambios fundamentais na Internet
        \item as aplicacións \textbf{reservan ancho de banda fim-a-fim}
        \item require software novo e complexo en hosts e roteadores
      \end{itemize}
    \end{column}
    \begin{column}{0.5\textwidth}
      \textbf{Filosofía \emph{laissez-faire}}
      \begin{itemize}
        \item sen cambios importantes
        \item máis ancho de banda cando faga falta
        \item distribución de contido, multicast na capa de aplicación
      \end{itemize}
      \textbf{Filosofía de servizos diferenciados (Diffserv)}
      \begin{itemize}
        \item menos cambios na infraestructura
        \item ofrece servizo de 1ª e 2ª clase
      \end{itemize}
    \end{column}
  \end{columns}
\end{frame}

\section{Clases de servizo e principios}

\begin{frame}{Proporcionar múltiples clases de servizo}
  \begin{itemize}
    \item ata agora: facer o mellor co servizo de ``mellor esforzo'' (unha única clase)
    \item alternativa: \textbf{múltiples clases de servizo}
      \begin{itemize}
        \item particionar o tráfico en clases
        \item a rede trata diferentes clases de forma diferente (analoxía: servizo VIP vs.\ normal)
      \end{itemize}
    \item granularidade: servizo \textbf{diferencial entre clases}, non entre conexións individuais
    \item historia: bits de \textbf{ToS} (IPv4) / \textbf{Flow Label} (IPv6)
  \end{itemize}
  \begin{center}
    \begin{tikzpicture}[font=\scriptsize]
      \node[draw,fill=aneliseAzul!13,rounded corners=2mm](h1){H1};
      \node[draw,fill=aneliseAzul!13,rounded corners=2mm,right=1cm of h1](h2){H2};
      \node[draw,fill=aneliseLaranja!25,rounded corners=1mm,below=1.2cm of h1](r1){R1};
      \node[draw,fill=aneliseLaranja!25,rounded corners=1mm,right=1.2cm of r1](r2){R2};
      \node[draw,fill=aneliseAzul!13,rounded corners=2mm,right=1cm of r2](h3){H3};
      \node[draw,fill=aneliseAzul!13,rounded corners=2mm,below=1.2cm of r2](h4){H4};
      \draw[thick] (h1) -- (r1);
      \draw[thick] (h2) -- (r1);
      \draw[thick] (r1) -- (r2);
      \draw[thick] (r2) -- (h3);
      \draw[thick] (r2) -- (h4);
      \node[below=1mm of r1,font=\tiny]{cola de saída R1: 1,5~Mbps};
    \end{tikzpicture}
  \end{center}
\end{frame}

\begin{frame}{Principio 1: marcado de paquete}
  \begin{exampleblock}{Escenario}
    Teléfono IP a 1~Mbps e FTP comparten enlace de 1,5~Mbps.
    \begin{itemize}
      \item as murmurios de FTP poden congestionar o roteador e causar perda de audio
      \item desexase dar \textbf{prioridade ao audio} sobre o FTP
    \end{itemize}
  \end{exampleblock}
  \begin{block}{Principio 1}
    \textbf{Marcado de paquete} necesario para que o roteador distinga entre clases; e nova \textbf{política de encamiñamento} para tratar os paquetes.
  \end{block}
\end{frame}

\begin{frame}{Principio 2: marcado e shaping na borda}
  \begin{itemize}
    \item e se as aplicacións se comportan mal (audio envía máis do declarado)?
    \item \textbf{Shaping}: forzar a adhesión da orixe ás alocacións de ancho de banda
  \end{itemize}
  \begin{center}
    \begin{tikzpicture}[font=\scriptsize]
      \node[draw,fill=aneliseAzul!13,rounded corners=2mm](tel){Teléfono 1~Mbps};
      \node[draw,fill=aneliseLaranja!25,rounded corners=1mm,right=2cm of tel](r1){R1};
      \node[draw,fill=aneliseLaranja!25,rounded corners=1mm,right=1.4cm of r1](r2){R2};
      \node[draw,fill=aneliseGris!15,rounded corners=2mm,below=1cm of r1](sh){shaper na borda};
      \draw[->,thick,aneliseLaranja] (tel) -- (r1);
      \draw[->,thick,aneliseLaranja] (r1) -- (r2);
      \draw[dashed,->,thick,aneliseVerde] (sh.north) -- (r1.south);
      \node[below=2mm of r2,font=\tiny]{enlace 1,5~Mbps};
    \end{tikzpicture}
  \end{center}
  \begin{block}{Principio 2}
    Marcado e shaping de paquete: \textbf{fornecer protección (illamento)} de unha clase fronte a outras.
  \end{block}
\end{frame}

\begin{frame}{Principio 3: uso eficiente dos recursos}
  \begin{itemize}
    \item alocar ancho de banda fixo (non compartible) a un fluxo: uso ineficaz se o fluxo non usa a súa alocación
  \end{itemize}
  \begin{center}
    \begin{tikzpicture}[font=\scriptsize]
      \node[draw,fill=aneliseAzul!13,rounded corners=2mm](tel){Teléfono 1~Mbps};
      \node[draw,fill=aneliseLaranja!25,rounded corners=1mm,right=2cm of tel](r1){R1};
      \node[draw,fill=aneliseLaranja!25,rounded corners=1mm,right=1.4cm of r1](r2){R2};
      \draw[thick,aneliseVerde] (r1.north) -- node[midway,above=1mm,font=\tiny]{enlace lóxico 1~Mbps} ([yshift=1cm]r1.north);
      \node[draw,fill=aneliseGris!15,rounded corners=1mm,below=2mm of r1](q1){enlace lóxico 0,5~Mbps};
      \draw[thick,aneliseLaranja] (tel) -- (r1);
      \draw[thick,aneliseLaranja] (r1) -- (r2);
      \node[below=2mm of r2,font=\tiny]{enlace 1,5~Mbps};
    \end{tikzpicture}
  \end{center}
  \begin{block}{Principio 3}
    Ao proporcionar illamento, é desexable usar os recursos da forma máis eficiente posible.
  \end{block}
\end{frame}

\begin{frame}{Principio 4: admisión de chamadas}
  \begin{columns}
    \begin{column}{0.55\textwidth}
      \textbf{Feito básico}: non se poden admitir demandas de tráfico máis alá da capacidade do enlace.
      \begin{center}
        \begin{tikzpicture}[font=\scriptsize]
          \node[draw,fill=aneliseAzul!13,rounded corners=2mm](t1){Tel 1~Mbps};
          \node[draw,fill=aneliseAzul!13,rounded corners=2mm,below=0.8cm of t1](t2){Tel 1~Mbps};
          \node[draw,fill=aneliseLaranja!25,rounded corners=1mm,right=2cm of t1](r1){R1};
          \node[draw,fill=aneliseLaranja!25,rounded corners=1mm,right=1.4cm of r1](r2){R2};
          \draw[thick,aneliseLaranja] (t1) -- (r1);
          \draw[thick,aneliseLaranja] (t2) -- (r1);
          \draw[thick,aneliseLaranja] (r1) -- (r2);
          \node[below=2mm of r2,font=\tiny]{enlace 1,5~Mbps};
        \end{tikzpicture}
      \end{center}
    \end{column}
    \begin{column}{0.45\textwidth}
      \textbf{Admisión de chamada}
      \begin{itemize}
        \item o fluxo declara as súas necesidades
        \item a rede pode bloquear a chamada (sinal de ocupado) se non pode atendela
      \end{itemize}
      \begin{alertblock}{Importante}
        A admisión de chamadas é a base da garantía de QoS fim-a-fim.
      \end{alertblock}
    \end{column}
  \end{columns}
\end{frame}

\section{Mecanismos de escalonamento e regulación}

\begin{frame}{Escalonamento: como se escolle o próximo paquete}
  \begin{columns}
    \begin{column}{0.5\textwidth}
      \textbf{FIFO (First In First Out)}
      \begin{itemize}
        \item enviar na orde de chegada á cola
        \item política de descarte: \emph{drop-tail}, prioridade, aleatorio
      \end{itemize}
      \textbf{Escalonamento prioritario}
      \begin{itemize}
        \item transmite da cola con maior prioridade
        \item múltiples clases con prioridades distintas
        \item a clase pode depender da marcación ou orixe/destino IP, portas\ldots
      \end{itemize}
    \end{column}
    \begin{column}{0.5\textwidth}
      \textbf{Varredura cíclica (round robin)}
      \begin{itemize}
        \item varre as colas das clases, atendendo unha de cada clase (se dispoñible)
      \end{itemize}
      \textbf{WFQ --- Weighted Fair Queueing}
      \begin{itemize}
        \item varredura cíclica xeneralizada
        \item cada clase recibe unha cantidade \textbf{ponderada} de servizo en cada ciclo
      \end{itemize}
      \begin{center}
        \begin{tikzpicture}[font=\scriptsize]
          \node[draw,fill=aneliseGris!15,rounded corners=1mm](cola){cola (espera)};
          \node[draw,fill=aneliseVerde!20,rounded corners=1mm,right=1.4cm of cola](link){enlace (servidor)};
          \node[draw,ellipse,fill=aneliseLaranja!25,above=0.6cm of link](s){escalonador};
          \draw[->,thick,aneliseLaranja] (s.south) -- (link.north);
          \node[above=0.5mm of cola,font=\tiny]{chegadas};
          \node[below=0.5mm of link,font=\tiny]{saídas};
        \end{tikzpicture}
      \end{center}
    \end{column}
  \end{columns}
\end{frame}

\begin{frame}{Mecanismos de regulación (shaping)}
  \begin{block}{Obxectivo}
    Limitar o tráfico para non exceder os parámetros declarados.
  \end{block}
  \textbf{Tres criterios comúns:}
  \begin{itemize}
    \item \textbf{taxa media} (longo prazo): cantos paquetes por unidade de tempo
      \begin{itemize}
        \item pregunta crucial: cal é o tamaño do intervalo? (100~pkt/s vs.\ 6000~pkt/min)
      \end{itemize}
    \item \textbf{taxa de pico}: ex.\ 6000~ppm media, 1500~ppm pico
    \item \textbf{tamaño da murmurio (máximo)}: nº máximo de paquetes consecutivos
  \end{itemize}
  \begin{block}{Balde de fichas (token bucket)}
    \begin{itemize}
      \item limita tamaño de murmurio e taxa media
      \item o balde mantén fichas; xéranse a taxa $r$ fichas/s (se o balde non está cheo)
      \item sobre intervalo $t$: paquetes admitidos $\leq r\cdot t + b$
    \end{itemize}
  \end{block}
\end{frame}

\begin{frame}{Token bucket + WFQ $\Rightarrow$ atraso acoutado}
  \begin{center}
    \begin{tikzpicture}[font=\scriptsize,
        c/.style={draw,rounded corners=2mm,fill=aneliseAzul!13,inner sep=1.2mm,text width=3cm,align=center},
      ]
      \node[c](tb){token bucket\\(taxa $r$, tamaño $b$)};
      \node[c,right=2.4cm of tb](wfq){WFQ\\(taxa por fluxo $R$)};
      \draw[->,thick,aneliseLaranja] (tb) -- node[midway,above=1mm,font=\tiny]{tráfico modelado} (wfq);
      \node[draw,fill=aneliseVerde!15,rounded corners=2mm,below=6mm of wfq](d){Atraso máx.\ $D_{\max}=b/R$};
      \draw[->,thick,aneliseVerde] (wfq.south) -- (d.north);
      \node[draw,fill=aneliseGris!12,rounded corners=2mm,below=6mm of tb](gar){Garantía de QoS};
      \draw[dashed,->] (tb.south) -- (gar.north);
    \end{tikzpicture}
  \end{center}
  \begin{itemize}
    \item o balde de fichas e WFQ combínanse para fornecer \textbf{límite superior garantido no atraso}
    \item isto é a base dos servizos garantidos de Intserv
  \end{itemize}
\end{frame}

\section{Diffserv e Intserv}

\begin{frame}{Servizos diferenciados da IETF (Diffserv)}
  \begin{itemize}
    \item queren clases de servizo ``cualitativas'': ``compórtase como un fío''
    \item distinción de servizo relativa: \textbf{Platinum, Gold, Silver}
    \item escalabilidade: \textbf{funcións simples no núcleo}, funcións complexas na borda
  \end{itemize}
  \begin{center}
    \begin{tikzpicture}[font=\scriptsize,
        c/.style={draw,rounded corners=2mm,fill=aneliseAzul!10,inner sep=1mm,text width=3cm,align=center},
      ]
      \node[c](borda){Roteador de borda\\marcado por fluxo\\(�perfil/fora de perfil)};
      \node[c,right=2.6cm of borda](nucleo1){Roteador de núcleo\\tráfico por clase\\buffering e escalonamento};
      \node[c,right=2.6cm of nucleo1](nucleo2){Roteador de núcleo\\(igual)};
      \draw[->,thick,aneliseLaranja] (borda) -- node[midway,above=1mm,font=\tiny]{preferencia a paquetes no perfil} (nucleo1);
      \draw[->,thick,aneliseLaranja] (nucleo1) -- (nucleo2);
    \end{tikzpicture}
  \end{center}
  \begin{itemize}
    \item 6 bits de \textbf{DSCP} (Differentiated Services Code Point) no campo ToS/Traffic Class
    \item determinan o \textbf{PHB} (Per-Hop Behavior) no roteador
  \end{itemize}
\end{frame}

\begin{frame}{Clasificación e condicionamento (Diffserv)}
  \begin{itemize}
    \item paquete marcado no campo \textbf{ToS} (IPv4) / \textbf{Traffic Class} (IPv6)
    \item 6 bits \textbf{DSCP} + 2 bits non usados (ECN)
    \item pode ser desexable limitar a taxa de inxección de tráfico dunha clase:
      \begin{itemize}
        \item o usuario declara \textbf{perfil de tráfico} (taxa, tamaño de murmurio)
        \item o tráfico mídese; móldase (shaping) se non está en conformidade
      \end{itemize}
  \end{itemize}
  \begin{block}{PHB (Per Hop Behavior)}
    \begin{itemize}
      \item PHB resulta en comportamento de repaso diferente nos roteadores
      \item non especifica cales mecanismos usar
      \item exemplos: \textbf{EF} (Expedited Forwarding), \textbf{AF} (Assured Forwarding, 4 clases $\times$ 3 particións de descarte)
    \end{itemize}
  \end{block}
\end{frame}

\begin{frame}{Servizos integrados da IETF (Intserv)}
  \begin{itemize}
    \item arquitectura para garantir QoS en redes IP para \textbf{sesións de aplicación individuais}
    \item \textbf{reserva de recursos}: os roteadores manteñen estado (tipo VC) de recursos alocados
    \item \textbf{admitir/negar} novas requirimentos de chamada
  \end{itemize}
  \begin{block}{Admisión de chamada}
    A sesión que chega debe:
    \begin{enumerate}
      \item declarar o seu requirimento QoS: \textbf{Rspec}
      \item caracterizar o tráfico: \textbf{Tspec}
      \item executar o protocolo de sinalización: \textbf{RSVP}
    \end{enumerate}
  \end{block}
  \begin{exampleblock}{Modelos de servizo}
    \textbf{Servizo garantido} (RFC 2212): chegada regulada por leaky bucket; atraso limitado matematicamente.
    \textbf{Carga controlada} (RFC 2211): ``QoS próxima á dunha rede non cargada''.
  \end{exampleblock}
\end{frame}

\begin{frame}{RSVP --- Resource Reservation Protocol}
  \begin{block}{Motivación}
    Encamiñamento sen conexión + mellor esforzo = sen sinalización no deseño inicial do IP.
    Novo requirimento: reservar recursos fim-a-fim para QoS en multimedia.
  \end{block}
  \begin{itemize}
    \item \textbf{RSVP} [RFC 2205]: ``permitir que os usuarios comuniquen requirimentos á rede de modo robusto e eficaz''
    \item antigo protocolo de sinalización: \textbf{ST-II} [RFC 1819]
  \end{itemize}
  \begin{block}{RSVP non\ldots}
    \begin{itemize}
      \item non especifica como reservar recursos (é un mecanismo de comunicación de necesidades)
      \item non determina as rutas (tarefa do roteamento); sinalización desacoplada do roteamento
      \item non interactúa co repaso de paquetes; separación de planos de control e datos
    \end{itemize}
  \end{block}
\end{frame}

\begin{frame}{RSVP: visión xeral da operación}
  \begin{columns}
    \begin{column}{0.55\textwidth}
      \begin{itemize}
        \item remitentes e receptores únense a un grupo multicast (fóra do RSVP)
        \item \textbf{sinalización do remitente á rede}:
          \begin{itemize}
            \item mensaxe \texttt{PATH}: fai coñecida a presenza do remitente
            \item \texttt{PATH teardown}: elimina o estado do camiño
          \end{itemize}
        \item \textbf{sinalización do receptor á rede}:
          \begin{itemize}
            \item mensaxe \texttt{RESV}: reserva recursos do(s) remitente(s)
            \item \texttt{RESV teardown}: elimina reservas
          \end{itemize}
        \item \textbf{sinalización da rede aos sistemas finais}:
          \begin{itemize}
            \item erro de camiño, erro de reserva
          \end{itemize}
      \end{itemize}
    \end{column}
    \begin{column}{0.45\textwidth}
      \textbf{Obxectivos de deseño}:
      \begin{enumerate}
        \item acomodar receptores heteroxéneos
        \item acomodar aplicacións con distintos requirimentos
        \item multicast como servizo de primeira clase
        \item aproveitar roteamento multicast/unicast
        \item overhead de protocolo lineal no nº de receptores
        \item deseño modular
      \end{enumerate}
    \end{column}
  \end{columns}
\end{frame}

\section{QoS en redes sen fío}

\begin{frame}{QoS en redes sen fío}
  \begin{itemize}
    \item idea: prover QoS en redes móbiles e ad hoc co \textbf{mínimo overhead posible}
    \item pódense usar os modelos de QoS das redes cableadas?
    \item dúas frontes: \textbf{QoS routing} e \textbf{QoS en Wi-Fi}
  \end{itemize}
  \begin{block}{QoS Routing}
    \begin{itemize}
      \item o roteamento é compoñente esencial: informa á fonte da banda e QoS dispoñible cara ao destino
      \item múltiples rutas poden ofrecerse con diferentes requirimentos de QoS
    \end{itemize}
  \end{block}
\end{frame}

\begin{frame}{QOLSR --- QoS para OLSR}
  \begin{itemize}
    \item OLSR é proactivo, estandarizado pola IETF (RFC 3626)
    \item QoS para OLSR proposto en 2003
    \item idea principal: usar \textbf{métricas de QoS na elección dos MPRs}
    \item as medidas de QoS obtéñense coas propias mensaxes \texttt{HELLO} --- \textbf{sen novas mensaxes de control}
  \end{itemize}
  \begin{block}{Heurísticas QOLSR}
    \begin{itemize}
      \item QOLSR1: selecciona o veciño de 1 salto con menor atraso como MPR
      \item QOLSR2: anticipa ademais o 2º salto (mínimo atraso 2-hop)
      \item QOLSR3: escolle MPR que alcance o maior número de nós descubertos
      \item non hai garantía de camiño óptimo respecto á métrica
    \end{itemize}
  \end{block}
  \begin{exampleblock}{Referencia}
    A.\ Munaretto, M.\ Fonseca, ``Routing and quality of service support for mobile ad hoc networks'', \emph{Computer Networks} 51:3142--3156, 2007.
  \end{exampleblock}
\end{frame}

\begin{frame}{QoS en Wi-Fi: anomalía de desempeño}
  \begin{block}{Performance Anomaly (802.11b)}
    \begin{itemize}
      \item nunha rede 802.11b, un único nodo lento (baixa SNR) reduce o throughput de todos os demais
      \item causa: CSMA/CA dá \textbf{igual oportunidade} de acceso independentemente da taxa
      \item o tempo do canal consúmese co nodo lento
    \end{itemize}
  \end{block}
  \begin{exampleblock}{Solución: Fair Time Sharing (FTS)}
    \begin{itemize}
      \item asignar \textbf{tempo de canal} proporcional á taxa do nodo (non slots iguais)
      \item A.\ Munaretto, M.\ Fonseca, ``Self-adapting algorithm to fair time sharing in wireless access networks'', \emph{Computers \& Electrical Engineering} 34:406--415, 2008
    \end{itemize}
  \end{exampleblock}
\end{frame}

\begin{frame}{802.11e --- QoS no MAC}
  \begin{columns}
    \begin{column}{0.5\textwidth}
      \textbf{EDCF / EDCA}
      \begin{itemize}
        \item 4 \textbf{Access Categories} (AC): Voice, Video, Best-Effort, Background
        \item cada AC ten seu \textbf{AIFS} e o seu \textbf{Backoff Counter}
        \item colisión virtual: a maior prioridade gaña, as outras incrementan CW
        \item \textbf{TXOP}: intervalo de transmisión sen contención
      \end{itemize}
    \end{column}
    \begin{column}{0.5\textwidth}
      \begin{center}
        \begin{tikzpicture}[font=\scriptsize,
            c/.style={draw,rounded corners=2mm,fill=aneliseAzul!13,inner sep=1mm,text width=3.3cm,align=center},
          ]
          \node[c,text width=3.4cm,align=center](ac0){AC0 Voice (AIFS curto)};
          \node[c,below=3mm of ac0,text width=3.4cm,align=center](ac1){AC1 Video};
          \node[c,below=3mm of ac1](ac2){AC2 Best-Effort};
          \node[c,below=3mm of ac2](ac3){AC3 Background};
          \node[draw,fill=aneliseLaranja!20,rounded corners=2mm,right=4mm of ac0,text width=2.6cm,align=center](vh){Virtual Collision Handler};
          \draw[->,thick,aneliseVerde] (ac0.east) to[bend left=10] (vh.west);
        \end{tikzpicture}
      \end{center}
      {\scriptsize Acceso inmediato se o medio idle $\geq$ AIFS[AC]$+$SlotTime.}
    \end{column}
  \end{columns}
  \begin{exampleblock}{Nota}
    IEEE 802.11e foi a base das garantías de QoS en Wi-Fi: sin ela, VoIP/vídeo sobre WLAN sofren.
  \end{exampleblock}
\end{frame}

\section{Conclusións}

\begin{frame}{Conclusións}
  \begin{alertblock}{Idea central}
    Prover garantías de QoS en redes sen fío é unha tarefa \textbf{complexa}!
  \end{alertblock}
  \begin{block}{Unha solución completa de QoS require}
    \begin{enumerate}
      \item un modelo apropiado de QoS
      \item un protocolo de sinalización de QoS
      \item un protocolo de roteamento con QoS
      \item un protocolo de acceso ao medio con QoS
      \item mecanismos adicionais: CAC, xestor de políticas, colas para control de conxestión\ldots
    \end{enumerate}
  \end{block}
  \begin{exampleblock}{Actualización 2026}
    As leccións de Intserv/Diffserv/RSVP viven en 5G \emph{network slicing} e en \emph{differentiation} Wi-Fi 6 (OFDMA + QoS).
  \end{exampleblock}
\end{frame}

\section{Bibliografía}

\begin{frame}{Bibliografía}
  \begin{itemize}
    \item \textbf{\kurose{}} --- Capítulo 7
  \end{itemize}
  \begin{block}{Adicional}
    \begin{itemize}
      \item Chakrabarti, Mishra, ``QoS issues in ad hoc wireless networks'', \emph{IEEE Communications Magazine} 39(2):142--148, 2001
      \item A.\ Munaretto, M.\ Fonseca, \emph{Computer Networks} 51:3142--3156, 2007 (QOLSR)
      \item A.\ Munaretto, M.\ Fonseca, \emph{Computers \& Electrical Engineering} 34:406--415, 2008 (FTS)
      \item Fonseca, Munaretto, Mendes, Chaouchi, ``A resource management framework for 802.11 wireless access networks'', \emph{Wireless Networks} 21:1891--1898, 2015
      \item RFC 2205 (RSVP), RFC 2212 (Guaranteed QoS), RFC 2211 (Controlled Load)
    \end{itemize}
  \end{block}
\end{frame}

\end{document}